\section{Conclusions (0.5 page)}
\label{sec:conclusions}
In conclusion, Minesweeper is cool.

\section{Future Work (0.5 page)}
\label{sec:future_work}
basically: do the project again, but bigger and more.
more minesweeper fields across a wider range of the 3bv spectrum \\
more minesweeper fields among the same 3bv value. \\

$RAC_{(i, m)}$ only separates minesweeper moves as far as the number of constraints $m$  and number of variables $i$ needed to consider for a pattern to be recognised (recall that Generalised Arc Consistency is given by $RAC_{(1,1)}$, whereas $RC_2$ and $RC_3$ use an unbounded i, effectively meaning only the number of constraints is relevant beyond GAC). Within the same number of constraints, e.g. $RAC_1$ therefore, different types of moves that use the same difficulty but are a different pattern are treated as the same by us.
Future work then should modify Relational Arc Consistency with a form of proof logging, to assert which kind of pattern was used to make a conclusion about a cell. this information can then be used to record what specific type of patterns a person knows and is able to use, and ultimately build a better understanding of their capability to play minesweeper.
\\ \\
use larger values for RAC. also, track guesses.
\\ \\
I build statistics on a minesweeper field, e.g. the typical difficulty for a cell,
by looking at many playthroughs and storing the difficulty of the cell as sampled by the players, at the moment they clicked the cell, during their playthroughs.
effectively, I'm using the player playthroughs as exploration of the board
instead of carrying out exploration myself
because I dont want to create an algorithm for exploring the board.

currently, I substitute trying every single possible way to solve the board and use player playthroughs instead.
I then build statistics from that.